\documentclass[11pt,a4paper]{article}

\usepackage{thga-db}

\renewcommand{\thgaKurs}{Databases}
\renewcommand{\thgaDatum}{Summer Term 2026}

\begin{document}
\thispagestyle{empty}

\thgaTitel{Term Project -- Documentation}%
{Vulnerability Tracker: a database-backed vulnerability management system}

\begin{infobox}
\textbf{Author:} Meriem Mornagui \\
\textbf{Module:} Introduction to Database Management Systems -- THGA Bochum \\
\textbf{Stack:} PostgreSQL -- FastAPI -- Docker Compose -- tkinter (\texttt{.deb}) \\
\textbf{Repository:} \url{https://github.com/Meriemmor/DBMS_10}
\end{infobox}

\tableofcontents
\newpage

% ============================================================
% ============================================================
% ============================================================
\section{Introduction and motivation}
% ============================================================

Modern IT systems may contain many assets and known security vulnerabilities.
Keeping track of which vulnerability affects which asset, who is responsible
for the affected system, and whether the issue has already been resolved can
quickly become difficult when this information is maintained manually.

The \textbf{Vulnerability Tracker} is a database-backed application for managing
security findings in a structured way. It stores \textbf{teams},
\textbf{assets}, \textbf{vulnerabilities}, and \textbf{findings}. A finding
connects a vulnerability to an affected asset and records its current status,
discovery date, due date, and, when applicable, resolution date.

The system uses PostgreSQL for persistent data storage and exposes the data
through a FastAPI REST API. A small tkinter desktop frontend communicates with
the API and allows a user to inspect and manage vulnerability findings without
accessing the database directly.

\begin{beispielbox}
\textbf{Usage scenario.} A security vulnerability is discovered on a server.
The user opens the Vulnerability Tracker, identifies the affected asset and
vulnerability, and creates a new finding. The responsible team can then track
the finding from \texttt{Open} to \texttt{In Progress} and finally to
\texttt{Resolved}. The summary view provides an overview of unresolved
findings grouped by severity.
\end{beispielbox}

% ============================================================
\section{Data model}
% ============================================================

The application manages four main entity types:
\textbf{teams}, \textbf{assets}, \textbf{vulnerabilities}, and
\textbf{findings}.

A team can be responsible for several assets. An asset can be affected by many
vulnerabilities, and one vulnerability can affect many assets. This creates an
\texttt{N:M} relationship between \textbf{asset} and
\textbf{vulnerability}. The relationship is resolved by the
\textbf{finding} entity, which also stores information about the state of the
vulnerability on the affected asset.

\begin{beispielbox}
\textbf{ER model (textual).}
\begin{itemize}[leftmargin=1.4em, topsep=2pt]
  \item \textbf{team} (1) \;--\; is responsible for \;--\; (N) \textbf{asset}
  \item \textbf{asset} (1) \;--\; has \;--\; (N) \textbf{finding}
  \item \textbf{vulnerability} (1) \;--\; appears in \;--\; (N) \textbf{finding}
\end{itemize}

Thus \textbf{asset} \texttt{N:M} \textbf{vulnerability}, resolved through
\textbf{finding}.
\end{beispielbox}

\subsection{Relational schema}

The relational schema derived from the ER model is:

\begin{itemize}[leftmargin=1.4em]
  \item \texttt{team}(
        \underline{id},
        name,
        contact\_email)

  \item \texttt{asset}(
        \underline{id},
        name,
        type,
        \emph{team\_id})

  \item \texttt{vulnerability}(
        \underline{id},
        cve\_id,
        severity,
        description)

  \item \texttt{finding}(
        \underline{id},
        \emph{asset\_id},
        \emph{vulnerability\_id},
        status,
        discovered\_on,
        due\_date,
        resolved\_on)
\end{itemize}

The foreign key \texttt{asset.team\_id} references \texttt{team.id}.
The table \texttt{finding} references both \texttt{asset} and
\texttt{vulnerability} and therefore implements the \texttt{N:M}
relationship.

\textbf{Normalisation.}
The schema is in \textbf{third normal form (3NF)}. Attributes describing teams
are stored only in \texttt{team}, asset-specific information is stored only in
\texttt{asset}, and vulnerability information is stored only in
\texttt{vulnerability}. The \texttt{finding} table stores only attributes that
describe the relationship between an asset and a vulnerability.

A unique constraint on \texttt{(asset\_id, vulnerability\_id)} prevents the
same vulnerability from being assigned to the same asset more than once.

% ============================================================
\section{Database (PostgreSQL)}
% ============================================================

The database is implemented with PostgreSQL. The schema and initial test data
are created automatically by the SQL script \texttt{db/init.sql}.

The database contains the four tables described in the data model:
\texttt{team}, \texttt{asset}, \texttt{vulnerability}, and
\texttt{finding}.

Several constraints are defined directly in PostgreSQL to ensure that invalid
data cannot be stored.

\begin{itemize}[leftmargin=1.4em]
  \item The field \texttt{vulnerability.cve\_id} is unique.
  \item The severity of a vulnerability must be one of
        \texttt{Critical}, \texttt{High}, \texttt{Medium}, or \texttt{Low}.
  \item The status of a finding must be one of
        \texttt{Open}, \texttt{In Progress}, or \texttt{Resolved}.
  \item The due date of a finding must not be earlier than its discovery date.
  \item A finding marked as \texttt{Resolved} must contain a resolution date.
        Unresolved findings must not contain one.
  \item The combination of \texttt{asset\_id} and
        \texttt{vulnerability\_id} is unique.
\end{itemize}

An excerpt from the SQL schema is shown below:

\begin{verbatim}
CREATE TABLE vulnerability (
    id SERIAL PRIMARY KEY,
    cve_id VARCHAR(50) NOT NULL UNIQUE,
    severity VARCHAR(20) NOT NULL
        CHECK (severity IN ('Critical', 'High', 'Medium', 'Low')),
    description TEXT
);

CREATE TABLE finding (
    id SERIAL PRIMARY KEY,
    asset_id INTEGER NOT NULL REFERENCES asset(id),
    vulnerability_id INTEGER NOT NULL REFERENCES vulnerability(id),
    status VARCHAR(20) NOT NULL
        CHECK (status IN ('Open', 'In Progress', 'Resolved')),
    discovered_on DATE NOT NULL,
    due_date DATE NOT NULL,
    resolved_on DATE,

    CHECK (due_date >= discovered_on),

    UNIQUE (asset_id, vulnerability_id)
);
\end{verbatim}

Foreign key constraints ensure referential integrity between the tables.
PostgreSQL therefore prevents findings from referencing assets or
vulnerabilities that do not exist.

The database is started as part of the Docker Compose backend. During the first
startup with an empty database volume, PostgreSQL executes the initialization
script and creates the required tables and example records.

% ============================================================
\section{REST API (FastAPI)}
% ============================================================

The backend API is implemented with FastAPI. It provides access to the
PostgreSQL database and is used by the desktop frontend to read and update
vulnerability data.

Read-only endpoints can be accessed without authentication. Endpoints that
modify data require an API key sent in the HTTP header
\texttt{X-API-Key}.

The implemented endpoints are:

\begin{center}
\begin{tabular}{|l|l|l|}
\hline
\textbf{Method} & \textbf{Endpoint} & \textbf{Purpose} \\
\hline
GET & \texttt{/findings} & List findings with optional filters \\
GET & \texttt{/findings/\{id\}} & Return one finding \\
GET & \texttt{/assets} & List assets \\
GET & \texttt{/teams} & List teams \\
GET & \texttt{/summary} & Show unresolved findings by severity \\
POST & \texttt{/findings} & Create a new finding \\
PATCH & \texttt{/findings/\{id\}} & Update an existing finding \\
\hline
\end{tabular}
\end{center}

The \texttt{/findings} endpoint combines data from multiple tables using SQL
joins. Each returned finding therefore contains information about the affected
asset, the responsible team, and the associated vulnerability.

The endpoint also supports optional filtering by status and severity. This
allows the frontend to request, for example, only open findings or only
critical vulnerabilities.

The \texttt{/summary} endpoint uses aggregation to count all unresolved
findings grouped by severity. The result contains counts for
\texttt{Critical}, \texttt{High}, \texttt{Medium}, and \texttt{Low}.

Write operations are protected by a simple API-key mechanism. The client must
send the correct key using the following request header:

\begin{verbatim}
X-API-Key: <api-key>
\end{verbatim}

If the API key is missing or invalid, the server rejects the request.

FastAPI also performs input validation using Pydantic models before data is
written to PostgreSQL. Invalid requests are therefore rejected before they can
modify the database.

% ============================================================
\section{Frontend (tkinter, packaged as .deb)}
% ============================================================

The desktop frontend is implemented in Python using \texttt{tkinter}. It
communicates exclusively with the FastAPI backend through HTTP requests and
does not access PostgreSQL directly.

When the application starts, a connection dialog is displayed. The user
provides the URL of the FastAPI server and the API key that is required for
write operations. After connecting, the main application window is opened.

The frontend provides the following main functionality:

\begin{itemize}[leftmargin=1.4em]
  \item Load and display vulnerability findings from the API.
  \item Display the summary of unresolved findings by severity.
  \item Refresh the displayed data.
  \item Create a new finding.
  \item Update the status of an existing finding.
\end{itemize}

Read operations are performed through the frontend's API helper function
\texttt{api\_get}. Operations that modify data use \texttt{api\_write} and
include the configured \texttt{X-API-Key} header.

For example, creating a finding opens a dialog in which the required finding
information can be entered. The frontend then sends the data to the
\texttt{POST /findings} endpoint. Updating the status of a finding uses the
\texttt{PATCH /findings/\{id\}} endpoint.

\subsection{Debian packaging}

The frontend is distributed as a Debian package so that it can be installed
as a normal application on the target Linux system.

PyInstaller is used to bundle the Python application, its required modules,
and the tkinter frontend into a standalone executable:

\begin{verbatim}
python3 -m PyInstaller \
  --onefile \
  --windowed \
  --hidden-import=tkinter \
  --name vulnerability-tracker \
  frontend/app.py
\end{verbatim}

The resulting executable is placed in
\texttt{frontend-deb/usr/bin/vulnerability-tracker}. The Debian package
metadata is stored in \texttt{frontend-deb/DEBIAN/control}.

The package is then built with:

\begin{verbatim}
dpkg-deb --build frontend-deb vulnerability-tracker_1.0_amd64.deb
\end{verbatim}

It can be installed on a compatible Debian-based system using:

\begin{verbatim}
sudo dpkg -i vulnerability-tracker_1.0_amd64.deb
\end{verbatim}

% ============================================================
\section{Operation and deployment}
% ============================================================

The backend consists of two services: the PostgreSQL database and the FastAPI
application. Both services are managed with Docker Compose.

The system can be started from the project directory with:

\begin{verbatim}
docker compose up -d --build
\end{verbatim}

Docker Compose starts the PostgreSQL container and the FastAPI backend
container. PostgreSQL stores the application data, while FastAPI provides the
REST interface used by the frontend.

The status of the containers can be checked with:

\begin{verbatim}
docker compose ps
\end{verbatim}

After startup, the API is available on port \texttt{8000}. The FastAPI
documentation can be accessed through the automatically generated Swagger UI
at:

\begin{verbatim}
http://<server-address>:8000/docs
\end{verbatim}

The desktop frontend is installed separately from the backend using the
generated Debian package. After installation, the application can be started
with:

\begin{verbatim}
vulnerability-tracker
\end{verbatim}

At startup, the user enters the backend URL and the API key in the connection
dialog. The frontend then communicates with the FastAPI service through HTTP.

For local development and testing, the complete backend has successfully been
started using Docker Compose and the API endpoints have been tested against the
PostgreSQL database.

The final deployment target is the lecture server. The backend containers and
the Debian frontend package will be installed and tested there before final
submission.

% ============================================================
\section{Reflection}
% ============================================================

The project demonstrates how the different components of a database-backed
application work together. The relational model provides a structured way to
represent teams, assets, vulnerabilities, and findings, while PostgreSQL
constraints help maintain data consistency.

One important part of the project was implementing the relationship between
assets and vulnerabilities through the \texttt{finding} table. This made it
possible to represent the many-to-many relationship while also storing
additional information such as status, discovery date, due date, and
resolution date.

The REST API provided practical experience with SQL joins, aggregation,
validation, and protected write operations. Separating the FastAPI backend
from the tkinter frontend also made the responsibilities of the different
application layers clearer.

Packaging the frontend was one of the more challenging parts of the project.
The application was bundled as a standalone executable with PyInstaller and
then packaged as a Debian package. This required ensuring that tkinter and the
necessary Python components were included correctly.

Overall, the project combines relational database design, SQL, REST APIs,
Docker, and desktop application development into one complete system.

% ============================================================
\section*{Sources and reuse}
% ============================================================

The project was developed based on the concepts and examples presented in the
DBMS course at THGA Bochum.

\begin{itemize}[leftmargin=1.4em]
  \item The relational database design, primary and foreign keys,
        normalisation, constraints, joins, and aggregation are based on
        concepts presented in the DBMS lectures and practical exercises.

  \item The PostgreSQL initialization approach and Docker Compose structure
        were adapted from examples used during the course and modified for the
        Vulnerability Tracker data model.

  \item The FastAPI backend structure and REST API concepts were adapted from
        the course examples and extended with the endpoints, validation,
        filtering, joins, aggregation, and API-key protection required by this
        project.

  \item The tkinter frontend and its communication with the REST API follow
        the frontend concepts presented in the course. The interface was
        adapted specifically for managing vulnerability findings.

  \item The frontend packaging process follows the course material on
        distributing Python desktop applications. PyInstaller is used to
        create the standalone executable, which is then packaged as a Debian
        package.
\end{itemize}

All database tables, application-specific API endpoints, validation rules,
frontend functionality, and the Vulnerability Tracker domain model were
implemented and adapted specifically for this term project.

\addcontentsline{toc}{section}{Sources and reuse}

\end{document}
